Теорема Ферма при $n = 3$ с использованием чисел Эйзенштейна (формулировка).
Метод спуска (свойствами $\lambda$ (задачи 1 -- 4) и леммой 
про сведение к уравнению упрощенного вида (задача 5) можно пользоваться б/д).

\textbf{Теорема}

Уравнение $x^3 + y^3 = z^3$ не имеет целочисленных решений.

\underline{Доказательство}

В предыдущем билете мы свели задачу к доказательству утверждения, что
$x^3 + y^3 = rz^3 \lambda^{3k}$ не может быть выполнено
ни для каких попарно взаимных простых 
$x, y, z \in \Z[\omega], r \in \Z{\omega}^*, k \in \N, 
\lambda \nmid xyz$.

Предположим, что это не так. Выберем такой набор чисел,
удовлетворяющих равенству, в котором число $k$ минимально.
Тогда (из прошлого билета) $k \geq 2$. Перепишем равенство
в виде 
$(x + y)(x + \omega y)(x + \omega^2 y) = rz^3 \lambda^{3k}$.

Заметим, что $(x + y, x + \omega y) = (x + y, \lambda y)$,
$x + y = (x + \omega y) + \lambda y$ и
$x + y = (x + \omega^2 y) + \lambda (1 + \omega) y$,
поэтому числа $x + y, x + \omega y, x + \omega^2 y$
имеют единственный общий делитель $\lambda$.
Значит, без ограничения общности их можно представить
в следующем виде:

\begin{equation}
\begin{gathered}
x + y = r_x x_0^3 \lambda, r_x \in \Z[\omega]^*, 
x_0 \in \Z[\omega] \\
x + \omega y = r_y y_0^3 \lambda, r_y \in \Z[\omega]^*,
y_0 \in \Z[\omega] \\
x + \omega^2 y = r_z z_0^3 \lambda^{3k - 2},
r_z \in \Z[\omega]^*, z_0 \in \Z[\omega]
\end{gathered}
\end{equation}

Но тогда, поскольку 
$(x + y) + \omega (x + \omega y) + 
\omega^2 (x + \omega^2 y) = 0$,
выполнено следующее:

\[
r_x x_0^3 + (r_y \omega) y_0^3 + 
(r_z \omega^2) z_0^3 \lambda^{3(k - 1)} = 0
\]

Тогда, как уже было доказано, $r_x = \pm r_y \omega$,
поэтому получен набор чисел,
удовлетворяющий равенству, с меньший показателем степени
у $\lambda$ -- противоречие.

